\documentclass[preprint, 3p,
authoryear]{elsarticle} %review=doublespace preprint=single 5p=2 column
%%% Begin My package additions %%%%%%%%%%%%%%%%%%%

\usepackage[hyphens]{url}

  \journal{Transport Geography?} % Sets Journal name

\usepackage{graphicx}
%%%%%%%%%%%%%%%% end my additions to header

\usepackage[T1]{fontenc}
\usepackage{lmodern}
\usepackage{amssymb,amsmath}
% TODO: Currently lineno needs to be loaded after amsmath because of conflict
% https://github.com/latex-lineno/lineno/issues/5
\usepackage{lineno} % add
\usepackage{ifxetex,ifluatex}
\usepackage{fixltx2e} % provides \textsubscript
% use upquote if available, for straight quotes in verbatim environments
\IfFileExists{upquote.sty}{\usepackage{upquote}}{}
\ifnum 0\ifxetex 1\fi\ifluatex 1\fi=0 % if pdftex
  \usepackage[utf8]{inputenc}
\else % if luatex or xelatex
  \usepackage{fontspec}
  \ifxetex
    \usepackage{xltxtra,xunicode}
  \fi
  \defaultfontfeatures{Mapping=tex-text,Scale=MatchLowercase}
  \newcommand{\euro}{€}
\fi
% use microtype if available
\IfFileExists{microtype.sty}{\usepackage{microtype}}{}
\usepackage[]{natbib}
\bibliographystyle{plainnat}

\ifxetex
  \usepackage[setpagesize=false, % page size defined by xetex
              unicode=false, % unicode breaks when used with xetex
              xetex]{hyperref}
\else
  \usepackage[unicode=true]{hyperref}
\fi
\hypersetup{breaklinks=true,
            bookmarks=true,
            pdfauthor={},
            pdftitle={Defining places for interchange: new tactial urbanist mobilityHubTools},
            colorlinks=false,
            urlcolor=blue,
            linkcolor=magenta,
            pdfborder={0 0 0}}

\setcounter{secnumdepth}{5}
% Pandoc toggle for numbering sections (defaults to be off)


% tightlist command for lists without linebreak
\providecommand{\tightlist}{%
  \setlength{\itemsep}{0pt}\setlength{\parskip}{0pt}}

% From pandoc table feature
\usepackage{longtable,booktabs,array}
\usepackage{calc} % for calculating minipage widths
% Correct order of tables after \paragraph or \subparagraph
\usepackage{etoolbox}
\makeatletter
\patchcmd\longtable{\par}{\if@noskipsec\mbox{}\fi\par}{}{}
\makeatother
% Allow footnotes in longtable head/foot
\IfFileExists{footnotehyper.sty}{\usepackage{footnotehyper}}{\usepackage{footnote}}
\makesavenoteenv{longtable}



\usepackage{subfig}
\usepackage{booktabs}
\usepackage{longtable}
\usepackage{array}
\usepackage{multirow}
\usepackage{wrapfig}
\usepackage{float}
\usepackage{colortbl}
\usepackage{pdflscape}
\usepackage{tabu}
\usepackage{threeparttable}
\usepackage{threeparttablex}
\usepackage[normalem]{ulem}
\usepackage{makecell}
\usepackage{xcolor}



\begin{document}


\begin{frontmatter}

  \title{Defining places for interchange: new tactial urbanist
mobilityHubTools}
    \author[]{James Reynolds%
  \corref{cor1}%
  \fnref{3}}
   \ead{julianj.reynolds@atu.ie} 
    \author[]{Noreen Armstrong%
  %
  \fnref{2}}
   \ead{????} 
    \author[]{Amaya Vega%
  %
  \fnref{2}}
   \ead{amaya.Vega@atu.ie} 
    \author[]{Brian McCann%
  %
  \fnref{3}}
   \ead{brian.mccann@atu.ie} 
    \author[]{Brian Caulfield%
  %
  \fnref{1}}
   \ead{brian.caulfield@tcd.ie} 
    \author[]{Agnieszka Stefaniec%
  %
  \fnref{4}}
   \ead{a.stefaniec@soton.ac.uk} 
    \author[]{Sarbast Molsem%
  %
  \fnref{1}}
   \ead{molsem@tcd.ie} 
      \cortext[cor1]{Corresponding author}
    \fntext[1]{Centre for Transport Research, Dept of Civil, Structural
and Environmental Engineering, Trinity College Dublin, Ireland}
    \fntext[2]{Atlantic Technological University, Galway City campus,
Ireland}
    \fntext[3]{Atlantic Technological University, Sligo campus, Ireland}
    \fntext[4]{Department of Decision Analytics and Risk at Southampton
Business School, University of Southampton, UK}
  
  \begin{abstract}
  
  \end{abstract}
  
 \end{frontmatter}

\section{Introduction}\label{introduction}

`Mobility hub' is a relatively new term for locations where a range of
transport modes, especially transit and shared micro-mobility modes, and
sometimes other services and facilities are offered. These are often
developed to provide a range of non-private-and-petrol-driven-car
alternatives to the surrounding communities and to interchanging
travellers and as part of efforts to shift towards more sustainable
travel behaviours and transport system outcomes. For many househouses
having nearby access to shared electric cars, bikes and cargo-bikes,
e-bikes and e-cargo bikes, or other vehicles might help in shifting
towards more walking, cycling and transit-based travel through providing
alternatives for those occassional trips that are otherwise difficult,
such as if travelling with luggage, longer distances or at a particular
time of day. Such households might be able to forgo purchase of a car or
an additional car, or hasten their transition to use of
electrically-powered modes that might have lower environmental impacts
without the need to purchase an e-car themselves.

However, what makes a location a `mobility hub', rather than just some
other place where there are a range of travel options in close
proximity, is unclear. Some mobility hubs, such as those in the Greater
Toronto and Hamilton Area (GTHA) in Canada, are defined as such by a
regional transport authority, with such status having implications for
landuse planning and other institutional decision-making
\cite{Salsberg:2010aa, Engel-Yan:2012aa}. Others are created afresh,
through some agency implementing facilities for a range of modes and/or
other offerings in close proximity to each other and branding as a
`mobility hub' or some other similar term. In some cases, it may be that
a range of offerings are already in place, especially where the private
sector has already implemented car-, scooter-, bike- or other -share
services, and converting somewhere into a `mobility hub' is within reach
by just giving the location a mobility-hub-related name in a
press-release or on a website, and perhaps some minimal physical
interventions such as signage and advertising.

Two problems, therefore, appear relevant: (1) that it is unclear what
minimum characteristics a location might need to become a `mobility
hub'; and (2) that there may not be sufficient tools available to
facilitate the transformation of such locations into `mobility hubs' by
either practitioners or, potentially through tactical urbanism and/or
grassroots activism, the general public. This paper reports case and
design-based research seeking to address these interlinked problems. The
research adopts the Double Diamond Design process by iterating through
alternative phases of diverging and coverging thinking related to
discovery, definition, development and delivery
\cite{UK-Design-Council:0aa}, with the discovery phase guided by
exploratory field visits to a range of case study mobility hubs. A new
package of software tools comprises the key research outcomes of the
delivery phase, with these providing a range of adaptable and
standardise mapping, branding and other outputs that might be used to
temporarily implement a `mobility hub', or as the basis for more
permanent signage, branding and other physical installations. The rest
of this paper is structured as follows: the next section outlines the
context for theis research, including existing knowledge related to
mobility hubs and tactical urbanism. The research methodology is then
described, followed by presentation of the results. This are then
discussed prior to a brief conclusion that includes the identification
of directions for future research.

\section{Research context}\label{research-context}

\subsection{Mobility hubs}\label{mobility-hubs}

Mobility hubs are sometimes described as part of an emerging approach
where shared and micro-mobility offerings might be linked, often with
transit, to provide a range of options that support more sustainable
travel. Early use of the term ``mobility hub'' occured in Toronto,
Canada, as part of landuse and transport planning undertaken by
Metrolinx and other governmental institutions across the GTHA and
surrounds \cite{Salsberg:2010aa, Engel-Yan:2012aa}. In ``The Big Move''
\cite{Metrolinx:2008aa} various transit stations and major interchanges
were defined as ``mobility hubs'' for the purposes of supporting
regional travel, modal interchange, and local connections including
through transitions to more transit supportive landuse patterns in
nearby areas.

The concept of having a range of transport modes, especially those that
are shared or that involve micro-mobility (e.g.~e-scooters, share
bikes), geographically collected into a `hub' may also have come about
from the plethora of independent private-sector transport offerings that
have recently emerged, including car-share and dockless bike/scooter
systems. Although driven largely by developments in the start-up,
technology and communications sectors, various governmental agencies
have sort to both regulate such offerings, but also make them act as
more than the sum of their parts, through defining dockless parking
locations, allocating dedicated car parking spaces and other
interventions, including through branding and informational signage that
provide geo-spatial context and collectivisation. As well, there have
been some emerginging problems with the various generations of dockless
bikes, e-scooters and e-bikes, including with respect to inappropriate
parking of vehicles (i.e.~blocking footpaths) and the fire safety and
environment impacts of recaharging models that rely on batteries being
collected in vans or small trucks and then charged, often in residential
or other properties without the appropriate facilities to deal with
electrical loads or battery fires \cite{Meigs:2023aa, Zhang:2022aa}.
Shared electric mobility hubs, such as those currently being trialled in
Ireland as part of the ROBUST project \cite{ROBUST:2025aa}, may provide
solutions to such problems by shifting the recharging of vehicles to be
at a hub itself, which includes the appropriate electrical connections
and equipment that, importantly for fires involving lithium batteries,
are outside and typically unenclosed.

Some challenges, however, for implementers and communities may relate to
the difficulty of developing mobility hubs, including the planning
processes and funding acquistion necessary to build new facilities. Of
course, there may be nothing especially new, other than the name, about
mobility hubs given that railway stations have existing for over a
century. Often containing platforms for boarding and alighting, waiting
areas and information, cafes, parking lots, interchange with other modes
such as buses, and many other facilities the 19th, 20th and current
iterations of railway stations would appear to have largely already
`solved' for the problem of providing support for mobility at a network
hub. However, railway stations tend to be located on land owned or
otherwise controlled by railway companies or agencies. As such,
facilities, signage, branding, information displays and most other
aspects of the station environment might be largely under the control of
a single institution, or at least a small number of closely-related
agencies. Mobility hubs, in contrast, might tend to be located on public
road reserves, within parking facilities or otherwise where there may be
less clear authority and power to make change without widespread
approvals or buy-in. These may also often require outside expertise
and/or involve top-down policy-making in which those making decisions
might be somewhat removed from the real mobility needs of individual
communities or may adopt `one size fits all' approaches that might lack
responsiveness to local context. Public policy analysis has long since
focused on such non-rational factors and how they might influence
decision-making, but these tend to be overlooked in transport policy
research \cite{Marsden:2017aa}, despite the applicability of
instititionalism, political approaches, incrementalism and other models
to transport \cite{Reynolds:2017ah}. Important within these are the
top-down, bottom-up and hybrid theories of implementation, which
emphasise the influence that street-level bureaucrats can wield to
obtain outcomes that respond to their local needs and context, including
sometimes outcomes that might be in opposition to central policy
directions. Further exploration of this topic may be beyond the scope of
this paper, but the extensive body of research related to the Advocacy
Coalition Framework (including
\cite{Sabatier:1986aa,Sabatier:1987aa,Sabatier:1988aa,Sabatier:1993aa,Weible:2009aa,Weible:2007aa})
provides insights and theory. Perhaps of more relevance, however, is the
field of tactical urbanism, which adopts many of the bottom-up and
grassroots approaches to defining and re-imaging public spaces,
especially those within road reserves and related to transport systems.

\subsection{Tactical urbanism}\label{tactical-urbanism}

Tactical urbanism refers to a range of approaches, often unsactioned,
for making changes to public environments. These often seek to question
what urban space or is for or the legitimacy of how it is allocated,
such as in the case of ``PARKing Day''. This started as lunchtime take
over of an on-street car parking space by an architecture firm, who
rolled out some green grass and a park bench, paid the parking meter,
and in doing so questioned the validity of using such spaces for storing
vehicles instead of as a paklet for people. ``PARKing Day'' is now an
annual event seeking, like the World Naked Bike Ride and Critical Mass,
to demonstrate the possibility of change and question the legitimacy of
the (car-dominated) status quo
\cite{Lydon:2012aa,Davidson:2013aa,Lydon:2015aa,Blumgart:2016ab, Ride:2009aa, Morhayim:2012aa, Longhurst:2015aa}.
Similar `pop-up' or protest-based approaches have helped legitimise
bicycle lanes, delegitimise car-use of bus lanes, and are starting to
becomem more mainstream (amongst engineers and planners) ways to help
build public and institutional support for new transit lanes
\cite{Fucoloro:2013aa,News:2019ab, TransitCenter:2018ab, Transportation-Studies:2019aa}.
These build on or have similarities to various earlier approaches
including street-reclaiming
\cite{Dobson:2023ae,Engwicht:1999aa,Fotel:2009aa,Institute:2015aa}, the
invention of the first `woornef' \cite{Guttenberg:1982aa} and Local Area
Traffic Management \cite{Damen:2017aa}, as well as being evident in the
make up of current directions towards `Movement and Place', `Complete
Streets' and Network Operation Planning approaches
\cite{Delbosc:2018aa}.

Despite the many names, variations and focii of these and other
invterventional and decision-making approaches, the call from
\cite{Marsden:2017aa} to ``pay greater attention to power\ldots{}''
retains relevance. Various authors have highlighted the challenges and
poptential hazards of participating in unsanctioned activities,
especially for minorities or others who might attract higher levels of
public surveillance or opposition, or attention from authorities, for
engagement in tactical urbanism or similar activities (INSERT CITATION
HERE ABOUT PARK BENCH). Intersections between illegality, illegitimacy
and the ill-enforced similarly exist in the context of street art and
its place on a contiunum that includes commissioned murals,
counter-culturism and the avent guard, graffiti, and all the way down to
plain old boring vandalism. \cite{Killen:2017aa,Killen:2021aa} explore
these topics in the context of limited operational and maintence impacts
to railways, specifically trains. This, and larger questions related to
`what is art' are again beyond the scope of this paper, but not those
that relate to who is allowed to post a public notice, define a
`mobility hub' or share knowledge about transport options, local
facilities and services, and connecting modes.

Engineers, planners, operational managers, communications teams and
others within transport agencies and governmental instutitions typically
have less need to engage in tactical-urbanist-type implementation,
typically because they have legitimate authority to make changes to the
built environment, at least within the scope of their remit. However,
pop-ups, trials, bottom-up and incremental approaches to implementation
may still be of use, or even necessary, for such actors seeking to
legitimise new, unusual or just controversial interventions, or if power
is vested in political representatives or otherwise.
\cite{Reynolds:2020aa,Reynolds:2021ab} defines these as `pragmatic
strategies for legitimation \emph{through} implementation' based on case
research of such approaches helping in the implementation of Bus Rapid
Transit in Curitiba and other forms of transit priority in Zurich,
Melbourne and Toronto. Other strategies include building legitimacy
before implementation using technical evaluation, vision-based transport
planning and/or public decision-making processes; or by avoiding impacts
on those who might seek to delegitimise change.

Such public processes and technical activities are readily apparent in
mobility hub implementation, as will be discussed in relation to some of
the individual case examples that are discussed later in this paper.
However, basic signage, linemarking or the public display of other
information would appear likely to be relatively easy for many
practitioners to implement within their own institutional setting. What
may be lacking, however, are tools and standarised approaches to
mobility hub definition and informational materials suitable for quick
and relatively easy adaptation to the local contex of an existing
mobility hub, or to `create' a mobility hub out of various closely
located modes, facilities and services.

Likewise, for activists, practitioners wishing to affect change outside
the remit of just the confines of their specific employment suituation,
or the general public there do not appear to be tools that might support
tactical urbanism efforts related to mobility hubs. This might be
particularly relevant for the `creation' of a mobility hub out of
already existing and geographically-proximate facilties, which might not
be immediately obvious to travellers but something well-known and
readily apparent for those with local knowledge of the geography and
community. Leaflets, posters for shop windows and community
noticeboards, or even the resources for a tactical urbanist to launch a
`guerilla marketing' campaign of posters, DIY signage and websites, and
likewise might be all that is needed to facilitiate bottom-up action and
empower communities to define their own hubs and advertise their
pre-existing mobility services and related facilities. What appears to
be lacking in the research literature, or in practice, are the tools to
produce the sorts of maps, analysis and interpretations that might
support the definition of a place as a hub and the communciation of its
existing and potential mobility and other offerings.

As discussed later in this paper, existing mapping, signage and similar
interventions in both the physical and virtual might tend to be produced
on a bespoke (ie. location by location) basis by specialist geographers
and graphic designers, and involve commercially-controlled datasets or
platforms. With the advent of crowd-sourcing, open data and the
widespread availability of open source software, however, there is the
possibilty of developing location-bespoke, but generalised, mapping and
similar tools and that might be useful for practitioners or others
wishing to `mobility hub-erise' some part of the urban environment. The
methods used in this research to firstly define what such tools might
do, and then to deliver a software package that enables others to more
easily develop their own mobility hub mapping and informational
materials, which might provide an initial version for future refinement,
are therefore descrbed in the following.

\section{Methodology}\label{methodology}

\cite{Coxon:2018aa,Coxon:2021ab} outline various industrial design
approaches and tools, as well as discussing their relevance and
potential applicability to transport systems. Notable amongst these is
the Double Diamond, which involves iteration across four phases,
alternating between diverging discovery and then converging definition
phases follwed by diverging development and convergance to a delivered
design output \cite{UK-Design-Council:0aa}. This study adopts this
approach as a framework to guide the research activity, which makes use
of case studies of existing mobility hubs for the initial discovery
phase.

Case studies can be especially effective for research into areas when
theory development is at an early stage and when there is a need to
explore phenomena within real-world contexts. Such is the situation
here, with the mobility hub concept still a relatively new development
in the transport planning and engineering fields, and a lack of
established practices or understandings. A challenge in using case-based
approaches, however, is the need to be seeking generalisable findings
despite researching at depth within the contexts of the individual case
or cases. This depth is possible through including only a small number
of cases in the study, typically up to five to ten but sometimes as few
as one, yet leads to the suggestion that researchers should explictly
respond to this ``duality criterion''
\cite{Benbasat:1987aa,Eisenhardt:1989aa,Dyer:1991aa,Cavaye:1996aa,Darke:1998aa,Meredith:1998aa,Stuart:2002aa,Voss:2002aa,Denscombe:2007aa,Eisenhardt:2007aa,Yin:2009ab,Ketokivi:2014aa,Yin:2014ab,Yin:2018ab}.

Here the need to be seeking generalisable findings is addressed directly
by the form taken by the research outputs: namely, the research produces
a software of package tools that might be used for mobility hubs in any
location, not just the locations selected as cases. More generally, the
case study (i.e.~field visits) component of this research is only part
of the diverging `discovery' and `development' phases, with additional
examples of mobility hubs examined through review of websites, online
imagery of street locations and other means. As such, while there may be
other approaches to information provision, mapping and branding at
mobility hubs not examined in this study, it is expected that the study
has enough `breadth' to offset the possible limitations to
generalisability that might come from examining only a small number of
cases a great `depth'.

\begin{figure}

{\centering \includegraphics[width=1\linewidth]{graphics/map_london} 

}

\caption{London hub locations. Source: authors based on Open Street Maps (OSM). OSM data and graphics copyright OSM contributors.}\label{fig:location_map_london}
\end{figure}

\begin{figure}

{\centering \includegraphics[width=1\linewidth]{graphics/map_birmingham} 

}

\caption{Hubs visited in Birmingham. Source: authors based on Open Street Maps (OSM). OSM data and graphics copyright OSM contributors.}\label{fig:location_map_birmingham}
\end{figure}

\begin{figure}

{\centering \includegraphics[width=1\linewidth]{graphics/map_bristol} 

}

\caption{Hubs visited in Bristol. Source: authors based on Open Street Maps (OSM). OSM data and graphics copyright OSM contributors.}\label{fig:location_map_bristol}
\end{figure}

\begin{figure}

{\centering \includegraphics[width=1\linewidth]{graphics/map_oxford} 

}

\caption{Hubs visited in Bristol. Source: authors based on Open Street Maps (OSM). OSM data and graphics copyright OSM contributors.}\label{fig:location_map_oxford}
\end{figure}

Figures \ref{fig:location_map_london} through
\ref{fig:location_map_oxford} show location of hubs visited in London,
Bristol, Brimingham and Oxford. Cases selection for this study was
somewhat constrained by circumstances and opportunity, at least for
those that were visited in person. The logisitics of and avaiable time
for travel, unfortunately, limited case selection to those that were
close to the researchers' home base in Ireland or that could otherwise
be visited as part of other travel. Hence, most of the visited sites
were in the south of England, including in the vicinity of London,
Oxford, Birmingham and Bristol (see Figures
\ref{fig:location_map_london} to \ref{fig:location_map_oxford}. However,
and in line with the recommendations of the methodological literature,
the cases selected for site visits include those that are leading and
particularly novel examples, those that are more representative, and
sufficient variation to provide for cross-case comparison. More
practically though, the actual conditions at each hub were only briefly
examined before each site visit, so as to retain some level of
randomness in the selection as far as the inclusion of good, average and
not as good examples of mobility hubs in the sampled sites. The location
of the mobility hubs themselves was sourced from (INSERT CITATION TO THE
MOBILITY HUB WEBSITE), which includes hundreds of sites across the UK
and Europe.

In the `definition' phase the findings from site visits and other
exploration of the cases are synthesised through cross-case comparison
so as to arrive at a design problem definition. The `diverging'
development phase is then guided by the standardise approaches to
package development for the R language, using \citet{wickham2023r} as a
key reference. While there are many other computing languages and
software tools that might have been selected instead, R is chosen here
because: it is open source and freely available, therefore supporting
the goals of this research that relate to developing an accessible suite
of tools for practitioners, advocates, tactical urbanists and others who
might not have immediate access to commercial software; it has an
extensive set of packages for geospatial analysis and mapping, including
\emph{sf} \citep{R-sf}, and the \emph{ggspatial} and \emph{osmdata}
\citep[\citet{R-osmdata}]{R-ggspatial} packages that facilitate access
to Open Street Maps (OSM) datasets and mapping tiles; and there is a
large community of users, the availability of beginner learning and
reference material\footnote{See for example INSERT THE MONASH COURSES
  HERE, AND THE O'REILLY ANIMAL BOOKS} and an online culture that
appears to support those with less experience with troubleshooting code
examples. This community of users is also responsible for the
development and maintenance of a wide range of R packages, including:
the \emph{tidyverse} \citep{tidyverse2019}; \emph{gtfstools}
\citep{R-gtfstools}; \emph{tidytransit} \citep{R-tidytransit}; and
INSERT OTHERS HERE, which were all relied upon to support the code
development aspects of this research.

\section{Results}\label{results}

\subsection{Discovery}\label{discovery}

\begin{figure}

{\centering \includegraphics[width=1\linewidth]{graphics/camden} 

}

\caption{Camden hubs (top three images) and other locations (rest), selected images. Source: authors.}\label{fig:camden_images}
\end{figure}

\begin{figure}

{\centering \includegraphics[width=1\linewidth]{graphics/south_woodford} 

}

\caption{South Woodford hub, selected images. Source: authors.}\label{fig:south_woodford_images}
\end{figure}

\begin{figure}

{\centering \includegraphics[width=1\linewidth]{graphics/wanstead} 

}

\caption{Wanstead hub, selected images. Source: authors.}\label{fig:wanstead_images}
\end{figure}

\begin{figure}

{\centering \includegraphics[width=1\linewidth]{graphics/andrew_road} 

}

\caption{Andrew Road hub, selected images. Source: authors.}\label{fig:andrew_road_images}
\end{figure}

\begin{figure}

{\centering \includegraphics[width=1\linewidth]{graphics/Bristol.001} 

}

\caption{Bristol hubs, selected images. Source: authors.}\label{fig:bristol_images}
\end{figure}

\begin{figure}

{\centering \includegraphics[width=1\linewidth]{graphics/Bristol.002} 

}

\caption{Bristol hubs information panels, selected images. Source: authors.}\label{fig:bristol_images_panels}
\end{figure}

\begin{figure}

{\centering \includegraphics[width=1\linewidth]{graphics/marina} 

}

\caption{Benson Marina hub, design plans and selected images. Source: authors, including plans from Oxfordshire County Council.}\label{fig:marina_images}
\end{figure}

The sites visited in Bristol share common design elements including:
elevated green cubes marking various hub components; seating within
parklet-style wooden decked areas; bicycle repair and maintenance
facilities; integrated USB A and C chargers (Figure
\ref{fig:bristol_images}, bottom-right); and extensive information
panels that include local and area-wide maps (Figure
\ref{fig:bristol_images_panels}. Those hubs located in Camden, likewise,
shared a common set of facilities: shared e-bike and e-cargo-bike,
secure bike parking and bike maintenance tools, and car club vehicles;
but not all of these were present at all of the hubs.As well and as
shown in Figure \ref{fig:camden_images}, there were many more
mobility-related offerings provided in the more central parts of London
on an individual basis, including: package pick-up; various bike hire
schemes both docked and dockless; and electric car charging. Hubs in
Brimingham likewise include directions, mapping and iconographic
representation of points of interest, facilities and other nearby
services that are not specifically `part' of a hub (Figure
\ref{fig:andrew_road_images}, bottom right), which highlights a
potential problem for implementers and designers: what services and
facilities to `include' as `part of the hub', versus those that are in
close proximity. In the case of the Transport for West Midlands hub at
Andrew Road, directional signage to and information about bus services
at the nearby Halesowen Bus Station (approx 50m walking distance) is
provided, but the hub has a different location-based name and branding
identity. For the proposed hub at Benson Marina, a cafe, the marina
itself and the Thames Path are only 30-40m from the nearest bus stop
(see Figure \ref{fig:marina_images}, purple circle, for location), but
while the design plans envisage a new footpath connection pedestrians
would still need to walk around a parking lot and upwards of 60 metres
to reach either. This is likely due to the ownership and access
restrictions, with the proposed works limited to the road reserve while
the marina, cafe and parking area appear to be part of a separate and
privately-held title.

Some of the hubs include linemarked areas for shared e-bike, e-scooter
or e-cargobike parking, bicycle parking and secure parking lockers,
although the offered facilities and services vary by location. The
proposed hub at Benson Marina, (Figure \ref{fig:marina_images}, provides
an example in that the design plans include an option include electric
car charging facilities and reserve parking spaces for shared cars.
Various potential operational and cross-institutional challenges,
however, were apparent. These include the challenges of cleaning and
maintenance, for example that vegetation planted in the parklets would
required regular attention and that planting boxes appeared to have
become a magnet for litter (Figure \ref{fig:bristol_images},
bottom-right).

A lack of shade, weather protection and drinking water was also evident
at most of the hubs, even though many of the parklet-type facilities
included steel framing that might have accommodated open-sided roofing
or similar. This might indicate an emphasis on form over function in
some of the element design, whereas one of the hubs in Bristol perhaps
provides an example of additional functionality, in a containerised
``Library of Things'' that appears to have been added later into the
vinicity of the hub, introduced onto top of areas linemarked and
sign-posted as for electric scooter hire (Figure
\ref{fig:location_map_bristol_images}, top-left and top-centre). This
may provide an indication of the possibilities of using modular
components to expand, repurpose or adjust mobility hubs as situations
change, but a need for signs, informational maps and other elements to
be able to be adjusted to match, perhaps with little to no involvement
of the original hub implementers. In the case of the ``Library of
Things'' addition at a hub in Bristol, the accommpanying green
informational sign (same design as Figure
\ref{fig:bristol_images_panel}) appeared to date from the original hub
implementation, and might be relatively challenging for practitioners,
location managers or others to keep `up to date' once the original
graphic designers or implementers have moved on\footnote{See also
  \cite{Reynolds:2017ah} discussion of `The Garbage Can Model'.}.

Overall, the site visits undertaken during the `Discovery' phase
highlight the following key issues: design elements; included facilities
and hub identity; maintenance; and the ability for a hub and supporting
informational material to be adjusted to changing circumstances. Design
elements relate to the desirability of functional and clearly related
hub components, and this appears interlinked with the lack of clarity
about what a mobility hub `is' and what types of nearby facilities
should be `part' of the hub in terms of branding and information. Hub
identity, likewise, appears to interlinked with few of the visited hubs
apparent have strong connections to named places. For example, despite
being close to the Halesowen bus station, a major interchange with eight
bus stands, the Andrew Road hub is named after a relatively less
important connector road that changes name to Highfield Lane a little
over 200m to the south west. Along with naming, the findings discussed
about generated through the site visits perhaps also highlight the
importance of maintenance, activity and updates to match changing
circumstances and give a hub a sense of `place'. Hubs in Bristol (Figure
\ref{fig:bristol_images}) perhaps highlight this, where litter,
non-mobility-related signage, space requisitioned for other purposes,
and some sites advertising dockless shared vehicles but without these
present in sufficient numbers (or at all at the time of visiting
Gainsborough Square) all provide indications of potential
post-implementation challenges related to mobility hubs.

\subsection{Definition}\label{definition}

This study seeks to address two research questions related to:
identifying locations with characteristics that may allow them to become
a `mobility hub'; and creating tools to support practitioners, tactical
urbanists or others in creating, operating and maintaining such
locations as mobility hubs. In the diverging Discovery phase discussed
above, the research used site visits to case study locations to build
understanding of mobility hubs and how they are defined in the real
world through informational signs, maps and other elements. This section
describes the subsequent and converiging Definition phase of the
research, in which a formal specification of the design problem to be
solved in the later Development and Delivery phases of the Double
Diamond process is developed.

As discussed above, findings from the sites visits suggest that there
are many different forms that a mobility hub might take, with a variety
of modes offered, services and facilities provided as part of the hub
itself, and services and facilities identified as being `nearby' as an
extension of the hub offer. This might suggest adopting an open-ended
and expansionist approach to mobility huib definition, given that the
designation might be applied to small sites with only a limited number
of modes, services and facilities just as much as to hubs at major
transport interchanges (e.g.~the Andrew Road hub's proximity to the
Halesowen Bus Station).

The consistency of branding at the Bristol hubs, especially with respect
to the iconographic green cubes drawing attention to hub elements, might
highlight the importance of proximity and line-of-sight in defining a
mobility hub. Although mapping at most of the visited hubs provided
information about services and facilities within up to around 15 minutes
walking distance, hub elements appeared to generally be located within
relatively confined areas, generally spanning less than 50 metres. The
Andrew Road hub again provides an example as an outlier, with the
entirety of the hub and the nearby Halesown Bus Station encompassed
within a radius of approximately 70 metres, and spanning slightly less
than 200 metres if traversed by foot. Hubs in London provide examples at
the other end of the spectrum, with the South Woodford hub taking up
approximately 15 metres of kerb space less than 100 metres from the
nearest tube station, while some in Camden and the ROBUST hubs
themselves use space equivalent to only a few car parking spaces.

Definitions of the facilities, services and modes necessary to become a
mobility hub appear likewise to be more likely to invovle open-ended and
inclusive approaches, rather than some sort of minimum checklist that
all sites must have. CoMoUK's pop-up mobility hub offering provide an
example, with the basic option consisting of a 12 square metre parklet
with bike repair facilties, and bike parking and/or bike share docking
\cite{CoMoUK:0aa}. In \emph{Mobility hub delivery models}
\cite{CoMoUK:2021aa} discuss six types of mobility hubs: ``large
interchanges/city hubs\ldots; transport corridor, smaller interchanges
or linking hubs; business park or ne housing development hubs;
suburb/mini-hubs; small market town/village hubs; and tourism hubs''
along with suggested components for each. Further types might included
the `shared' and `electric' forms of mobility hub, as in the ROBUST
sites, wherein the offered modes are available for hire and electrically
powered respectively. It is unclear, however, whether there is a need to
differentiate between electric hubs (or `e-hubs') where the
electrification involves on-site recharging, as at ROBUST and for many
e-cars, and hubs where electric vehicles are dockless or otherwise
`recharged' through the replacement of battery packs with those powered
up elsewhere. \cite{Meigs:2023aa} and \cite{Zhang:2022aa} discuss some
of the problems associated with this first form of electrification,
while \cite{Reynolds:0ab} highlights the constraints of access to power
infrastructure and issues related to electrical connections for the
latter, which might limit location choice and extend implementation
timelines. Guidance from \cite{CoMoUK:0ab} suggests a timeline of 12 to
15 months to step through consultation, design, stakeholder engagement,
tendering and construction, although this appears to relate to either
unpowered mobility hubs or the former type of electricification. This
guidance also suggests that budgets in the vicinity of 250k to 650k
pounds \cite{CoMoUK:0ac}; and that ``mobility hubs work best when part
of a wider sustainable transport network that provides active travel,
public transport and shared transport options'' \cite{CoMoUK:0aa}.

The design problem, therefore, appears to relate to determining how to
create mobility hubs that: incorporate existing modes, facilities and
services; minimise cost and timeframes, such that practitioners and
tactical urbanists might build legitimacy through implementation (c.f.
\cite{Reynolds:2020aa,Reynolds:2021ab}) using minimal resources,
incremental approaches that allow expansion one-hub-at-a-time, and that
are scaleable; and allow for rapid iteration, revision or adjustment to
account for changing circumstances. There are limits, however, to what
can be achieved in the context of this research study, suggesting that
the focus here should be on building tools that allow for mobility hub
creation without the sorts of `hard' engineering works that would
necessitate design, tendering and construction of new facilities or
infrastructure. Rather, this research defines the design problem to be
solved as to create software tools that might facilitate: the
identification of locations that have the potential to become `mobility
hubs'; the creation of user information in both virtual and physical
environments (e.g.~websites, posters, maps); maintenance and revision of
developed hubs to support changing circumstances (i.e.~hub element
removal or addition) and community ownership (i.e.~open-source, freely
available and without the need for extensive specialist tools or
skills);and rapid development to support bottom-up, tactical urbanist
and iterative approaches.

\subsection{Development}\label{development}

As shown in Figure \ref{fig:bristol_images} and
\ref{fig:bristol_images_panels}, the various elements the Bristol hubs
are linked through visual cues provided by the elevated green cube
markers and matching iconography on the information panel and local
maps. Even though some components of some of the hubs may not be
immediately visible from the main panel and/or separated by up to around
30-50 metres, this iconography appears to provide consistent, logical
and clear local wayfinding around the hub facilities themselves. As
discussed in relation to the previous `discovery' phase, however, some
facilities (e.g.~toilets) are shown on parts of the main informational
panels (the top and the map), but not marked in the physical world
(i.e.~with the green cubes), while other elements such as cafes,
groceries shops and similar are only shown on the maps.

Figure \ref{fig:andrew_road_images}, top-left and bottom-left, show
various hub elements (e.g.~electric car charges, bike docks) on more
temporary signage, including: (top-left) replaceable poster of the type
typical at transit stops; and (bottom-left) a banner sign printed on
flexible materials. Notably, the Andrew Road hub is billed as a ``Local
Travel Point pilot'' and at the time of the site visit there was no
information provided about the, at the time incomplete, electric car
charging facilities. This might imply that information at Andrew Road is
to be updated as the pilot progress, and the mobility offer or other
elements change. This contrasts to the permanency of the mapping and
information in Bristol (Figure \ref{fig:bristol_images_panels}), which
includes iconography and descriptive text detailing the hub offerings
and providing information about nearby services. Such examples from the
cases might highlight the value to practitioners, activists and others
of easily updatable mapping and informational products

\subsection{Recreating Bristol
elements}\label{recreating-bristol-elements}

A starting point for Development is to recreate the Bristol style of hub
imformation, mapping and similar elements. Here this is achieved through
the use of Open Streets Map (OSM) data sets and mapping tiles, as well
as various software packages within the R programming language. Icons
are drawn from \cite{Pousson:2022aa}

\subsubsection{Adjusting and
generalising}\label{adjusting-and-generalising}

LAT AND LON branding? As per that Collingwood Brew Pub shirt I had.

Identifying `a hub'

Identifying the `facitiles and services' to be included as `part' of
each hub, be they

Connected naming, branding etc.

Website, posters, banners, leaflets and other informational products
that might be generated quickly by practitioners, adovocates or others,
and easily maintained and updated.

Dated to show how current the information is. Version control may also
be important.

More generic parts of the puzzle - maintenance schedules, inspection
logs, management resources and other items that might help a mobility
hub `product owner' facilitate operations, continuance and (maybe even)
growth.

\subsubsection{Create Richmond maps}\label{create-richmond-maps}

\subsection{Delivery}\label{delivery}

\section{Discussion}\label{discussion}

\section{Conclusions}\label{conclusions}

\begin{longtable}[]{@{}llllllll@{}}
\caption{Author contribution matrix, as per CRediT (Contribution Roles
Taxonomy}\tabularnewline
\toprule\noalign{}
Component & NA & BC & BM & SM & JR & AS & AV \\
\midrule\noalign{}
\endfirsthead
\toprule\noalign{}
Component & NA & BC & BM & SM & JR & AS & AV \\
\midrule\noalign{}
\endhead
\bottomrule\noalign{}
\endlastfoot
Conceptualization & & & & & X & & \\
Data Curation & & & & & X & & \\
Formal Analysis & & & & & X & & \\
Funding Acquisition & & X & & & & X & \\
Investigation & & & & & X & & \\
Methodology & & & & & X & & \\
Project Administration & & & & X & & & \\
Resources & & & & & & & \\
Software & & & & & X & & \\
Supervision & & & X & & & & X \\
Validation & & & & & X & & \\
Visualization & & & & & X & & \\
Writing -- Original Draft Preparation & & & & & X & & \\
Writing -- Review \& Editing & & & & & & & \\
\end{longtable}

\section*{References}\label{references}
\addcontentsline{toc}{section}{References}

\renewcommand\refname{Appendicies}
\bibliography{References.bib, packages.bib}


\end{document}
